\documentclass{article}
\begin{document}
\section*{Comparing independent block evaluation with evaluation of the whole path}
For every ordered pair of natural lengths, a fixed split native grade span and fixed ordinary category-valued terminal presheaf, the independently staged evaluation is the functor obtained by joining the two strictly matched blocks, applying the old whole-path evaluator, and transporting along the arithmetic equality of the output presheaves. This equality holds on the entire weighted total category, not only on identity-base arrows. The base, grades and payload use object universe $u$ and hom universe $\max(u,v)$. Matching at the internal base is strict, on arrows as well as objects. This is not an assertion of equality for an unrelated weak pullback presentation.
\par Encoding a moving packed arrow retains the full ordered list of primitive arrows, with endpoints, plus the moving terminal payload arrow; splitting one packed layer by the genuine square equivalence has a left inverse. Head recursion proves injectivity at all lengths. Independent snoc recursion on the native positive powers identifies their last grade, including transport along the earlier terminal square. The actual native block projections concatenate these words. Suffix-first staging produces that same trace, proving equality by injectivity before any use of packed multiplication. The arithmetic transport between the double and whole iterates never permutes grades.
\par Restriction to the strict source fibre and descent across ordered charts give the natural isomorphism between the \#211 staged and \#210 full-path functors. Whiskering and pasting yield the square with actual weighted native flattening and genuine packed multiplication. The old lifting interface recovers precisely its pre-multiplication transformation. Both zero cases are included without reducing positive suffixes to an assumed square.
\par The proof does not assert that the forgotten cut itself can be reconstructed from one multiplied path. In the ordered-cut probe $[4,6,8]$, cuts after one and two grades produce the same full path but have different prefix endpoints. The moving-arrow probe uses Boolean endomaps on a one-object base, varying the internal grade while keeping exterior, terminal and payload observations equal. Terminal reindexing distinguishes opposite composites with words $[4,9,2,6]$ and $[8,4,9]$. The source probe compares genuine tagged naturality on noninvertible arrows with terminal contexts 6 and 10, retaining different noncommuting payload words.
\par Remaining: joint variation of this tagged compositor with cleavage-based source change and changes of payload or split span, coherence under repeated cutting and units, global monad comparison and strict contextad-algebra recovery. Issue \#105 stays open.
\end{document}
